\section{Laplace transform and transfer functions}

\subsection{Ctrl+F keywords}

\subsubsection{Laplace, differential equation, transfer function, poles, zeros, poler, nulpunkter}
\[
\text{ODE, characteristic equation, zero initial conditions, proper, strictly proper}
\]
Use this section when an exam problem gives a differential equation or asks for poles, zeros, stability, or \(G(s)\).

\subsection{Laplace rules}

\subsubsection{Unilateral Laplace transform}
\[
X(s)=\mathcal{L}\{x(t)\}=\int_{0^-}^{\infty}x(t)e^{-st}\,dt
\]
The transform converts linear differential equations into algebraic equations.

\subsubsection{Derivative property with initial conditions}
\[
\mathcal{L}\{x^{(n)}(t)\}=s^nX(s)-\sum_{k=0}^{n-1}s^{n-1-k}x^{(k)}(0^-)
\]
For transfer functions, set all initial conditions to zero.

\subsubsection{Common transform pairs}
\[
\begin{array}{c|c}
x(t) & X(s)\\ \hline
\delta(t) & 1\\
u(t) & 1/s\\
t u(t) & 1/s^2\\
e^{-at}u(t) & 1/(s+a)\\
\sin(\omega t)u(t) & \omega/(s^2+\omega^2)\\
\cos(\omega t)u(t) & s/(s^2+\omega^2)
\end{array}
\]
These pairs are enough for the recurring transfer-function and final-value tasks.

\subsection{Core transfer-function formulas}

\subsubsection{Transfer function definition}
\[
G(s)=\frac{Y(s)}{U(s)}\bigg|_{\text{zero initial conditions}}
\]
A transfer function is only the input-output relation for the selected input and output.

\subsubsection{Differential equation to transfer function}
\[
a_n y^{(n)}+\cdots+a_1\dot{y}+a_0y
=b_m u^{(m)}+\cdots+b_1\dot{u}+b_0u
\]
\[
G(s)=\frac{Y(s)}{U(s)}
=\frac{b_ms^m+\cdots+b_1s+b_0}{a_ns^n+\cdots+a_1s+a_0}
\]
Denominator roots are poles. Numerator roots are zeros.

\subsubsection{Poles, zeros, gain form}
\[
G(s)=K\frac{\prod_i(s-z_i)}{\prod_k(s-p_k)}
\]
Poles decide natural modes and stability. Zeros shape response and Bode phase.

\subsubsection{Properness and implementability}
\[
\deg N(s)\leq \deg D(s)\quad\text{proper},\qquad
\deg N(s)<\deg D(s)\quad\text{strictly proper}
\]
Controllers, filters, and feed-forward blocks must be proper for direct implementation.

\subsection{Exam recipe: ODE to transfer function and poles}

\textbf{Use when:} the problem gives a differential equation and asks for \(G(s)\), poles, stability, or a multiple-choice transfer function.

\textbf{Steps:}
\begin{enumerate}
\item Set initial conditions to zero.
\item Replace \(y^{(k)}\) by \(s^kY(s)\) and \(u^{(k)}\) by \(s^kU(s)\).
\item Move all \(Y\)-terms to one side and \(U\)-terms to the other.
\item Form \(G(s)=Y(s)/U(s)\).
\item Factor or inspect the denominator for poles.
\end{enumerate}

\textbf{Fast checks:} highest derivative of \(y\) sets denominator order; signs in the ODE carry into the polynomial.

\textbf{Common traps:} keeping nonzero initial-condition terms inside a transfer function; confusing poles of \(G\) with zeros of \(G\).

\subsection{Exam recipe: characteristic equation from feedback}

\textbf{Use when:} the problem asks for closed-loop poles or stability of a feedback system.

\textbf{Steps:}
\begin{enumerate}
\item Find the loop transfer \(L(s)\).
\item Write the characteristic equation \(1+L(s)=0\) for negative feedback.
\item Clear denominators to get a polynomial.
\item Test pole locations or use margin/Nyquist methods when the polynomial is not requested.
\end{enumerate}

\textbf{Fast checks:} numerator zeros of \(T=L/(1+L)\) do not decide closed-loop stability unless they cancel unstable poles, which should not be assumed.
